\section{Centering, coercivity, and the fluctuation spectrum}\label{sec:03-gas-mean}

The finite contour formula expresses the cylinder partition function as
an oscillatory sum of particle integrals. We first identify its exact
mean density. Subtracting this density leaves an interaction whose real
part is positive and therefore controls both the particle counts and
excursions beyond a logarithmic wall. We then compare the interior
interaction with a reference ensemble whose normalization and linear
statistics can be evaluated. The energy estimates have two uses. The untruncated real energy estimate controls every valid particle count; the sharper imaginary trial fields identify the tilt cost in the fixed-fugacity exponent. Both use the same exact compensated-sector identity, Lemma~\ref{lem:03-compensated-sector}. The determinant asymptotic in Theorem~\ref{thm:03-determinant} supplies the remaining fluctuation contribution. Lemmas~\ref{lem:03-core} and~\ref{lem:03-schur} then give the finite-dimensional approximation needed to realize that determinant by moving actual particle coordinates.

\subsection{The exact mean density}

Throughout this section we specialize Theorem~\ref{thm:02-contour-sectors} to \(N=2m,\ s=b=m,\ \varepsilon=1\), all real site offsets zero. Fourier transforms in the real coordinate use \(e^{-iku}\) and inverse measure \(dk/(2\pi)\). A particle is specified by its real coordinate \(u\) and its ray label; the imaginary part of its logarithmic coordinate is determined by that label. An interaction matrix kernel acts by convolution; symmetric means \(H^{JI}(-u)=H^{IJ}(u)\), and forms with complex arguments are continued \emph{bilinearly}. Write
\[
 \gamma=\pi/(3\beta),\quad L=(\log N)/\gamma,\quad
 p=2\beta,\quad \kappa=(3\beta)^2 .
\]
Labels are \(I=(\sigma,J)\), where \(\sigma\in\{+1,-1\}\) and \(J=0,F,P\) denote the main, far-single and far-pair species. Their elementary shifts are, respectively, \(0\), \(11/4\), and \((5/2,7/2)\), as in Theorem~\ref{thm:02-contour-sectors}. Their multiplicities \(w_I\) are \(1,1,2\); a pair still has only one real integration variable. Write \(H\) for minus the sign-normalized particle-particle log interaction, \(V\) for the log (without negation) of the sign-normalized site-particle interaction (particle row, site column). Differences in these kernels take the first real coordinate minus the second, e.g. the exponentiated factors use \(H^{I I'}(u-u')\) with minus sign for labels \(I,I'\). More explicitly use linear combinations of
\[
 K_\theta(u)=-\log[2(\sin(\theta/2)\cosh(u/2)-i\cos(\theta/2)\sinh(u/2))],\qquad
 K_0(u)=-\log|2\sinh(u/2)|
\]
reducing angles to \([0,2\pi)\). Thus \(H\) uses the exponents of the log-sinh factors in \(E\) (fused for pair labels), and \(V\) the negative exponents in the site-particle product; cross terms in reverse order transpose and reflect the kernel. The exponents at angle zero in \(H\) are \(2\delta_{IJ}\) by the shift lists; there are no angle-zero terms of \(V\). These conventions agree with the sign normalization of the contour-sector identity.

Modulo distributions at zero frequency the transforms of \(K_\theta\) are
\[
 \pi e^{(\pi-\theta)k}/(k\sinh(\pi k)),\qquad
 \widehat K_0=\pi\cosh(\pi k)/(k\sinh(\pi k)).
\]
Indeed \(K_\theta''=\tfrac14/\sinh^2((u+i\theta)/2)\); shifting by a period and taking the double pole residue gives \(-\pi k e^{(\pi-\theta)k}/\sinh(\pi k)\). The second formula follows by even averaging as \(\theta\downarrow0\). Affine terms in comparisons below are fixed by tails: real parts \(-|u|/2+O(e^{-|u|})\), imaginary parts odd with limits given by \(\tfrac12(\pi-\theta)\operatorname{sign}u\) for \(0<\theta<2\pi\). For the main (near) labels,
\[
 \widehat H_{nn}(k)=\frac{2\pi c_*(k)}{k\sinh(\pi k)}
 \begin{pmatrix}\cosh(3\beta k)& e^{pk}\\ e^{-pk}&\cosh(3\beta k)\end{pmatrix},
 \qquad c_*(k)=2\cosh(\beta k)-1.
\]
Define \(\widehat\varrho=1/c_*\). Shifting by a period and taking simple residues gives
\[
 \varrho(u)=\frac{e^{-\gamma|u|}-e^{-5\gamma|u|}}{\sqrt3\,\beta(1-e^{-6\gamma|u|})}.
\]
The value at zero is its removable continuation. This is a positive,
even, analytic probability density, and \(\varrho(u)\asymp
e^{-\gamma|u|}\). The lattice activity is denoted by
\(\rho\).

Let \(\mu\) have density \(m\varrho\) on each main label and vanish
on the four far labels. We use \((f,Hg)\) for the sum of the
corresponding double integrals over labels and real coordinates.

\begin{lemma}[Exact centering]\label{lem:03-centering}\label{E:centering}
The particle--site kernel satisfies \(V=H_{\cdot n}*\varrho\).
Combining the site--site factors with
\(\exp((\mu,H\mu)/2)\) gives
\[
 C_*^N\prod_{i<j}H_*(a_i-a_j)^2,\qquad
 C_*=\frac{H_*(0)}{\sqrt{B_s(0)}}
     =\frac{(1+d)\gamma}{4d\cos(\beta/2)},
\]
with the fixed signs prescribed by the contour-sector formula.
Here the logarithmic site coordinates \(a_i\) are the balanced confluent sites: \(m\) copies of \(0\) and \(m\) copies of \(-i\beta\). The function \(H_*\) is the removable analytic function used in Theorem~\ref{thm:02-denominator}.
\end{lemma}
\begin{proof}
On the main rows the identity follows from the displayed Fourier symbols. For an elementary far shift \(j\) in group \(\sigma\), divide the \(H\)-row versus near columns (own group first) by \(2\pi c_*/(k\sinh\pi k)\); it is
\((\cosh(\beta k)e^{\sigma(4-j)\beta k},e^{\sigma(p-j\beta)k})\), as follows using \(2<j<5\). These entries give \(V\) after convolution for far shifts as well. Equality at nonzero frequencies suffices (matching real tails, with zero limiting difference, and bounded odd imaginary parts).

For a site pair, the sign-normalized log of \(B_s\) or \(B_c\) plus the corresponding entry of \(\varrho*H_{nn}*\varrho\) at real log difference equals the normalized log of \(H_*^2\) at the \emph{site} log difference (with imaginary shift). For the same-group formula the symbol divided by \(2\pi/(k\sinh\pi k)\) on the left is
\(-\cosh(2\beta k)-\cosh(\beta k)+\cosh(3\beta k)/c_*(k)\).
On the right it is \(-2\cosh(\beta k)+\cosh(4\beta k)-\sinh(\pi k)\tanh(3\beta k/2)\), the same expression. To check the second expression, factor \(\cosh u+d\) and \(\sinh(u/2)\) in \(H_*\), and use the scaled \(K_\pi-K_0\) identity
\[
 \widehat{\log|\tanh(\gamma u/2)|}(k)
       =-\frac{\pi}{k}\tanh\frac{\pi k}{2\gamma}.
\] Cross entries acquire the factor \(e^{-\beta k}\) (extra sinh terms cancel since \(\sinh(3\beta k)/c_*=\sinh(2\beta k)+\sinh(\beta k)\)); \(H_*\) is analytic and nonzero through this shift, so its log acquires the same shift in the transform of the second derivative. Again real tails have matching slopes and zero constant, and imaginary parts are bounded odd. In particular the exponentiated half self term of \(\varrho*H_{nn}*\varrho\) is
\[
 C_*=\frac{H_*(0)}{\sqrt{B_s(0)}}=\frac{(1+d)\gamma}{4d\cos(\beta/2)} .
\]
In \(B_s(0)\) the argument is the log difference. Thus all one-particle terms can be centered exactly using \(H\mu\); site pairs together with the mean factor \(\exp((\mu,H\mu)/2)\) yield \(\prod_{i<j}H_*(a_i-a_j)^2\) times \(C_*^N\) (up to fixed signs).

\end{proof}

\subsection{Positivity across the six species}

Exact centering is useful only if the remaining energy controls every
species, including far particles. We prove that control on arbitrary
real fluctuations; no sign condition on the main components is needed.
For a real fluctuation \(x=(x_I)\) of zero \emph{weighted total} mass, write
\[
 d_x=\sum_I w_I x_I,\quad D=\partial^{-1} d_x,\quad r=\sum_{I:\sigma=+1} w_I x_I,\quad
 B(x)=\tfrac12(x,\operatorname{Re}H\,x)
\]
with primitives starting at \(-\infty\). Initially take smooth rapidly decreasing densities.

\begin{lemma}[Coercivity of the real interaction]\label{lem:03-coercivity}\label{E:coercivity}
Suppose that the components of \(x\) are initially smooth rapidly decreasing real densities and that \(\int d_x=0\). The tail \(-w_Iw_J|u|\) contributes exactly \(\|D\|_2^2\); residual real kernels are integrable. We have
\[
 B(x)\asymp\|D\|_2^2+\sum_I\|x_I\|_{-1/2}^2,\qquad
 \|x_I\|_{-1/2}^2=\int(1+k^2)^{-1/2}|\widehat x_I|^2\,dk/(2\pi).
\]
The form extends accordingly and controls \(\|D\|_{1/2}^2=\int (1+k^2)^{1/2}|\widehat D|^2\,dk/(2\pi)\) also.

\end{lemma}
\begin{proof}
Divide the real-kernel transform by \(\pi/(k\sinh\pi k)\); the matrix in group blocks has diagonals \(S\), off-diagonals \(T\). With \(h_j=\cosh(j\beta k/4)\), upper triangles listed by rows, they are
\[
\begin{split}
 S:\ &2(h_{16}-h_{12}+h_8),\ h_{13}-h_1+h_3+2h_5-h_9,\ h_{14}+h_2+2h_6;\\
     &2(h_{16}-h_{12}+h_8),\ h_{13}+h_{15}+h_5+h_7;\\
     &2(h_{12}+h_{16}+h_4+h_8).\\
 T:\ &2(h_{12}+h_4-h_8),\ 2(h_1-h_3+h_7),\ 2(h_{10}+h_2);\\
     &2h_{10},\ 2(h_{11}+h_9);\\
     &4(h_{12}+h_8).
\end{split}
\]
These follow by summing log-factor exponents with angles modulo \(2\pi\). Expand in \(v=\cosh(\beta k/4)-1\). At zero \(S=T=2w w^{\mathsf t}\), \(w=(1,1,2)^{\mathsf t}\). All coefficient matrices of \(S\pm T\) of positive degrees through 16 are strictly positive definite after division by the corresponding coefficient of \(h_{16}\). Use the ratios
\[
 R_i(l)=i^2\prod_{h=1}^{l-1}(i^2-h^2)\big/\big(16^2\prod_{h=1}^{l-1}(16^2-h^2)\big)
\]
by the cosh polynomial differential equation \((2v+v^2)h_i''+(1+v)h_i'=i^2 h_i\) as polynomials in \(v\). For clarity the first three normalized \(S,T\) triangles are:
\[
\begin{array}{c|lll}
 l&S_{11},S_{12},S_{13}&S_{22},S_{23}&S_{33}\\\hline
1 &11/8,73/128,17/16&11/8,117/64&15/4\\
2 &203/136,483/1360,849/1360&203/136,1703/1360&939/340\\
3 &685/408,5849/22848,23/51&685/408,3163/3264&1133/476
\end{array}
\]
\[
\begin{array}{c|lll}
 l&T_{11},T_{12},T_{13}&T_{22},T_{23}&T_{33}\\\hline
1 &3/4,41/128,13/16&25/32,101/64&13/4\\
2 &35/68,19/272,413/1360&165/544,175/272&513/340\\
3 &9/28,293/22848,55/476&55/476,2035/7616&155/204
\end{array}
\]
Both sum and difference have positive successive leading minors (they exceed \(1/2,1/4,1/10\) respectively by substitution). For degrees \(l\ge4\), using \(0\le R_i(l)\le R_i(4)\) gives absolute off-diagonal bounds \(0.21,0.39,0.90\) and diagonal lower bounds \(1.60,1.75,1.59\), enough by diagonal dominance. This coefficient calculation gives the asserted norm comparison as follows. Near zero, the constant rank-one term contributes \(\|D\|_2^2\), and the strictly positive degree-one coefficients control the remaining directions at constant scale. On each compact frequency interval avoiding zero, positivity and continuity give uniform upper and lower bounds. At large \(|k|\), the \(h_{16}=\cosh(4\beta k)\) term is diagonal and dominates all other entries. Since \(\pi=4\beta\), multiplication by \(\pi/(k\sinh\pi k)\) gives the scale \(|k|^{-1}\) there. These three ranges prove the two-sided Sobolev bound. Finally \(\widehat d_x=ik\widehat D\), and \(d_x=\sum_Iw_Ix_I\), so the same bound controls the stated \(H^{1/2}\)-norm of \(D\).

\end{proof}

The full norm comparison controls counts. The leading low-frequency
terms will also identify the optimal imaginary tilt. In the following
calculation, Fourier transforms are implicit on the displayed
variables. At small \(k\) the pointwise quadratic symbol of \(B\), using variables \(D,r,x_{\rm far}\) and suppressing hats for these variables, is
\[
 |D|^2+Q(r,x_{\rm far})+O(|k|)(|D|^2+|r|^2+|x_{\rm far}|^2).
\]
Here \(Q\) is a constant positive definite real form and \(Q(r,0)=(\kappa-p^2)|r|^2\). Indeed insert \(d_x=ikD\) in the matrices just given. With rows versus near columns, \(\operatorname{Im}H\) has a term proportional to \(\operatorname{sign}u\); the remaining odd kernel \(J_o\) has multiplier bounded by \(C\min(|k|,(1+|k|)^{-1})\) (integrable with moments, piecewise smooth of bounded variation). For the main rows the sign coefficients form the antisymmetric matrix with upper entry \(p\). A far row differs from \(w_I\) times its corresponding main row; the additional coefficients against \(+,-\) are sums over shifts \(j\) of
\[
 ((4-j)\beta,-j\beta)\ \ (\sigma=+),\qquad
 (j\beta,(j-4)\beta)\ \ (\sigma=-).
\]
In particular the difference of the two additional coefficients is \(\pi w_I\), with the same sign on both groups of far rays. Write \(g_I\) for the second (summed) additional coefficient.

\subsection{Interior and explicit particles}

We retain the collective fluctuations of particles near their mean and
estimate the remaining particles separately. Split the main contours
at \(\pm L\). Keep as explicit particles \(\zeta\) all main ones outside the wall and all far particles everywhere, regarded as positive counting measures by label (unordered, with the factorial measures of the sector formula). Write \(n_e\) for their total number of variables, \(f=\sum_{I:\rm far}w_I\int\zeta_I\), and let
\[
 m'_\sigma=m-\sigma M-\sum_{I:\sigma_I=\sigma}w_I\int\zeta_I .
\]
The remaining interior variables are main-label particles with totals \(m'_\sigma\). Splitting including factorials is exact by symmetry. The physical sector has nonnegative integer \(m'_\sigma\). In
variational estimates we also use signed near densities
\(\alpha_\sigma\) on \([-L,L]\) with these same prescribed totals,
and we allow \(M\) to be complex. This enlargement concerns trial
densities only; it does not change the integer dimensions of a particle
integral.

We regularize only the interaction \emph{between} interior and explicit particles (denoted by subscripts \(n,e\)), as follows. Fix \(h>0\) small; it stays fixed through \(N\to\infty\). Take a positive convolution probability \(\psi=\psi^{(h)}\) with transform the average of
\(\cosh((\pi-ht)k)/\cosh(\pi k)\) for \(0<t<1\), with density proportional to \(\exp(-1/\sqrt t)\). These kernels in real space before averaging are
\[
 \frac{\cosh(u/2)\sin(ht/2)}{2\pi(\sinh^2(u/2)+\sin^2(ht/2))}
\]
by Fourier inversion (elementary sech integral by residues and imaginary translation). Thus \(\psi\) is even, bounded by \(C/h\), with tail bounded by \(C h e^{-|u|/2}\) at \(|u|\ge1\), and \(\int u^2\psi\le C h\). Its transform is bounded above and below by \(C\exp(-c(h|k|)^{1/3})\) and \(c\exp(-C(h|k|)^{1/3})\) respectively, by optimizing \(ht|k|+1/\sqrt t\) (for the lower bound integrate a comparable interval); \(1-\widehat\psi\le Ch k^2\).

Put \(S_0=2K_0\). Then \(S_0\ge S_0*\psi\ge S_0*\psi*\psi\) on the real line, since convolution once replaces \(K_0\) by averaged \(\operatorname{Re}K_{ht}\), by the transform formulas with matching tails. Let \(\chi\) be an even smooth cutoff equal to 1 on \([-1/2,1/2]\), zero off \([-1,1]\), \(0\le\chi\le1\), with derivatives bounded by \(C^{l+1}(l!)^2\). Such cutoffs follow by integrating/rescaling bumps formed from \(\exp(-1/x)\) on \(x>0\) zero-extended, by Cauchy estimates on circles of radius \(c x\). Subtract from \(H_{ne}\), only on same-label entries (thus involving an exterior main particle), the real nonnegative kernel
\[
 R_h(u)=\chi(u)(S_0-S_0*\psi)(u).
\]
Call the resulting interaction \(\widetilde H_{ne}\), extending \(R_h\) by zero on other entries. The interaction \(\widetilde H_{ne}\) is smooth Gevrey (some finite order, for fixed \(h\)) with at most linear growth and bounded higher derivatives of each positive order. Here and below Gevrey bounds follow equivalently for bounded decaying terms from stretched exponential Fourier estimates; for example differentiating under a bound \(C\exp(-c|k|^\tau)\) gives \(C^{l+1}(l!)^{1/\tau}\) bounds, and multiplication of Gevrey functions preserves a finite order by Leibniz. For \(S_0*\psi\) one can first take two derivatives in the Fourier estimates.

For smooth trial densities define the complex energy
\[
 \mathcal E(\alpha,\zeta)=\tfrac12(\alpha-\mu,H_{nn}(\alpha-\mu))
 +(\alpha-\mu,\widetilde H_{ne}\zeta)+\tfrac12(\zeta,H\zeta)_{\rm off}.
\]
The last term omits all individual explicit self terms. One may fix generic distinct explicit real positions (null sets of coincidences never needed). Trial bounds will apply for example to smooth densities, or integrable densities of finite log energy with sufficient decay, by approximation. Point configurations have logarithmic self singularities, so this
formula must be compensated before it is used in a particle integral.
The following reference kernel makes that compensation explicit. Let \(u\mapsto t\) be an increasing bijection \([-L,L]\to[0,\pi]\) smooth with strictly positive derivative in the open interval; set \(p_L=dt/\pi\) as a probability there in \(u\) (we will specify the map). Use
\[
 H_{\rm ref}(t,t')=-2\log|2(\cos t-\cos t')|
\]
separately on each main group, zero cross. Its \(p_L\)-mean vanishes by circle reflection (mean of \(\log|2\sin(t/2)|\) is zero by its Fourier series). Use this kernel for independent reference integrals with measures
\[
 \frac{1}{\nu_\sigma!}\prod_{i=1}^{\nu_\sigma}\frac{2\sin t_i\,dt_i}{2\pi}
 \prod_{i<j}|2(\cos t_i-\cos t_j)|^2 .
\]
Under $x=\cos t$ this is the Legendre orthogonal-polynomial ensemble.
Its normalization by polynomial norms uses Andr\'eief's identity for
the integral of a product of two determinants; see
\cite[equations (1.3), (3.1), and Section 5]{TracyWidom1998}.
The same identity will give
the tilted factorial densities in the fluctuation calculation.
For interior configurations \(X_\nu\) with these variables, put \(q_\sigma=X_{\nu,\sigma}-\nu_\sigma p_L\). The smooth (near-diagonal compensated) observable uses
\[
 \Phi_M(q)=\mathcal E(m'p_L+q,\zeta)-\tfrac12(q,H_{\rm ref}q).
\]
This notation denotes a single smooth-kernel pairing, not a difference
of two infinite empirical self energies. More explicitly, put
\(a=m'p_L-\mu\) and
\[
 K=H_{nn}-H_{\rm ref},\qquad
 b=H_{nn}a+\widetilde H_{ne}\zeta.
\]
The reference kernel is understood in the physical coordinates on
each main component. Since its \(p_L\)-mean is zero,
\[
 \Phi_M(q)=\mathcal E(m'p_L,\zeta)+(b,q)+\tfrac12(q,Kq).
\]
All diagonal values in the last pairing are the finite continuous
values of \(K\); the coordinate construction below proves the needed
smoothness. This is also the meaning of \(\Phi_M\) when
\(q=X_\nu-\nu p_L\) is a centered empirical distribution.
\begin{lemma}[Compensated sector representation]\label{lem:03-compensated-sector}
Fix an explicit configuration \(\zeta\) with generic distinct real coordinates and a physical integer imbalance \(M\) for which \(m'_\sigma\ge0\). After division of the right-hand side of \eqref{eq:02-contour-sectors} by
\((\gamma/2)^N\prod_{i<j}H_*(a_i-a_j)^2\), its contribution at \(M,\zeta\) is
\[
 c_\zeta\,e^{i(\eta+\pi f)M}
 \int A(X_{m'},\zeta)\exp\{-\Phi_M(X_{m'}-m'p_L)\},
 \qquad A(X,\zeta)=\exp\{-(X,R_h\zeta)\}.
\]
The integral includes the two \emph{unnormalized} reference measures. The explicit integration measures and their species factorials are retained outside it. The prefactor \(c_\zeta\) is independent of \(M\), satisfies \(|c_\zeta|\le\exp(C_hn_e)\), and uses the fixed sign convention of Lemma~\ref{lem:02-sector-signs}. Moreover \(0\le A\le1\). Sectors with a negative \(m'_\sigma\) contribute zero.
\end{lemma}
\begin{proof}
Restoring the omitted self terms of \(H_{nn}-H_{\rm ref}\) contributes per interior variable
\[
 (c/d)\,2\sin t\,t'(u),
\]
since \(E_s(u)\sim(d/c)^2u^2\). Combined with \(du/(2\pi)\), this is \(c/d\) times the one-particle reference measure. The scalar cancellation is
\[
 (2d)C_*(c/d)=\gamma/2.
\]
Replacing \(\widetilde H\) by \(H\) gives exactly the additional factors \(A\) and \(\exp((\mu,R_h\zeta))\). The latter is at most \(\exp(C_hn_e)\): the cutoff can meet \(\mu\) only within distance one of an exterior main particle, where \(|u|\ge L-1\), and its logarithmic singularity is integrable. Missing \(c/d\) factors and far activities each cost a fixed constant per explicit variable. Finally the sector sign calculation leaves precisely the phase \(e^{i(\eta+\pi f)M}\), with all remaining signs independent of \(M\) at fixed \(\zeta\). This proves both the identity and the prefactor bound.
\end{proof}

The identity is algebraic on distinct interior configurations, hence holds almost everywhere in the reference integrals. The coordinate choice below gives the smooth reflected extension of \(K\) needed for distributional estimates at the walls. For later interpolation we also \emph{define} the same compensated integrand with complex \(M,m'\) and independent nonnegative integer reference dimensions \(\nu\), always setting \(q=X_\nu-\nu p_L\). This extension keeps \(q\) centered at the actual integer reference dimension; only \(\nu=m'(M)\) gives the physical identity of the lemma.

We now compare this regularized energy with the positive real form.
The loss from omitting the explicit self interactions will be linear
in the number of explicit variables.

\begin{lemma}[Uniform smoothing inequality]\label{lem:03-smoothing}
For the trial energy defined above, set
\[
 x=\operatorname{Re}\alpha-\mu+\psi*\zeta,\quad z=\operatorname{Im}\alpha
\]
with all species included for \(x\), only near entries nonzero for \(z\). These are weighted-total neutral. Then
\begin{equation}
 \operatorname{Re}\mathcal E(\alpha,\zeta)
 \ \ge\ B(x)-B(z)-(x,\operatorname{Im}H\,z)
 -s(h)(B(x)+B(z)+1)-C_h n_e,\qquad s(h)\longrightarrow0.
 \label{eq:G-1}
\end{equation}

\end{lemma}
\begin{proof}
At nonzero frequency,
\[
 |\widehat{\widetilde H}_{ne}-\widehat H_{ne}\widehat\psi|
 \le C h^{c_0}(1+|k|)^{-1}\widehat\psi(k)
\]
for some \(c_0>0\), with no residual distributions at zero. Indeed nonsingular terms have transform \(O(e^{-c|k|}(1+k^{-2}))\), multiplied by \(1-\widehat\psi\). For the singular \(S_0\) term the difference in real space is \((1-\chi)(S_0-S_0*\psi)\) with derivatives in \(L^1\) bounded by \(O(h^2)C^{l+1}(l!)^2\) by even averaging of imaginary shifts off the singularity, and the cutoff bounds. This gives \(O(h^2)\exp(-c\sqrt{|k|})\) in frequency, sufficient by the lower bound for \(\widehat\psi\). Thus the mixed-kernel error is bounded by
\(C h^{c_0}\|\alpha-\mu\|_{-1/2}\|\psi*\zeta\|_{-1/2}\).
For explicit-explicit pairs only the real kernel matters. Every smooth real summand (apart from diagonal \(S_0\)) changes under \(\psi*\psi\) averaging by \(O(h)e^{-c|u|}\), using exponentially decaying second derivatives and the tail and second moment of \(\psi\). The diagonal singularity can only decrease. And the missing same-variable diagonals in the smoothed model cost \(O_h(n_e)\).

To absorb these errors we need a quadratic count bound that is valid even when the interior trial densities are signed. We obtain it by testing far labels separately and testing main labels only outside the wall, where those trial densities vanish. Partition the explicit coordinates into unit-scale cells and let \(n_j\) denote their variable counts. We claim
\begin{equation}
 \sum n_j^2\le C(1+|\log h|)(B(x)+1).
 \label{eq:G-2}
\end{equation}
For far labels use translates of a nonnegative fixed bump covering a unit interval, to test separately each \(x_I=\psi*\zeta_I\) positive. For outer main counts on either side of the wall use bumps supported \emph{outside} the wall, adapted to cells starting from the wall (the grid may change by bounded overlaps). For the first cell use a bump rising from zero at the wall to 1 within distance \(h\), plateauing to distance 2, then cutting off. This captures at least a constant smeared mass from each unsmeared point of the first cell, by the displayed real kernels. On all these outer tests \(\operatorname{Re}\alpha\) does not intervene and \(\mu\) has cell masses summably bounded. The \(H^{1/2}\) norms squared of the bumps are \(O(1+|\log h|)\) even at the wall, by the kernel seminorm \(\int|b(u)-b(v)|^2|u-v|^{-2}du\,dv\) (equivalence by Plancherel); for the transition of length \(h\), the translation differences in \(L^2\) squared cost \(O(\min(l,l^2/h))\) for translations by \(0<l<1\). Bounded overlap gives the same squared-norm bound times \(\sum a_j^2\) for linear combinations (sum differences pointwise has bounded multiplicity). Duality proves \eqref{eq:G-2}.

Also \(\|\psi*\zeta\|_{-1/2}^2\le C(1+|\log h|)\sum n_j^2\). Indeed the spatial kernel for multiplier \((1+k^2)^{-1/2}\) is bounded by \(C(1+|\log|u||)e^{-c|u|}\), by writing the symbol as a gamma-integral of Gaussians. Convolved twice with \(\psi\) it is bounded by \(C(1+|\log h|)e^{-c'|u|}\) using the density and tail bounds. This and \eqref{eq:G-2} control the mixed error just found (substitute \(\alpha-\mu=x+i z-\psi*\zeta\)), and the \(O(h)e^{-c|u|}\) explicit pair errors, proving \eqref{eq:G-1}.

\end{proof}

\begin{corollary}[Control of explicit particles]\label{cor:03-counts}\label{E:explicit-counts}
Constants may depend on the fixed smoothing parameter \(h\). One has
\begin{equation}
 |\operatorname{Re}M|^2+n_e^2\le C_h L(B(x)+1),\qquad
 \ell_e:=\sum_{v\in\zeta} (|v|-L)_+\le C(B(x)+1).
 \label{eq:G-3}
\end{equation}
\end{corollary}
\begin{proof}
Outside the wall \(\operatorname{Re}\alpha\) is absent. For \(t>L\), the total count beyond \(t\) is at most twice its positive smeared cumulative there, hence its square integrates on \([L,\infty)\) to \(O(B(x)+1)\) by \(\|D\|^2\) and the \(\mu\)-tail. Same on the other side. The cumulative counts are nonnegative integers, so their squares dominate the counts themselves. Integrating gives the distance sum bound. It also bounds exterior counts away from the wall; \eqref{eq:G-2} controls the remaining cells near the wall and the far-particle cells inside it.

Total \emph{smeared} mass outside the wall is at most \(C_h(1+\sqrt{B(x)})\) also, by the cumulative test outside distance one from the wall and \eqref{eq:G-2} convolved by \(\psi\) for the first cell. Finally \(\int r=-\operatorname{Re}M\); test on the wall region with a smooth plateau extending outside by distance one, cost \(O(\sqrt L)\) in \(H^{1/2}\), and bound the remainder by the positive tail masses. In particular, retaining a fraction of \(B(x)\) in a trial lower bound gives exponential damping \(\exp(-c_h\ell_e)\), sufficient for convergence and tail truncation of the explicit absolute integrals. Summing their factorial measures with factors \(\exp(C_hn_e)\) costs at most \(\exp(O_h(L))\): far variables have effective volume \(O(L)\), and exterior main variables have effective volume \(O(1)\). This is the coarse cost; part 4 of Proposition~\ref{prop:03-trial-bounds} will improve it for the configurations needed by the exponent.

\end{proof}

\subsection{The energy gain from an imaginary tilt}

The cylinder parameter enters a sector through
\(e^{i(\eta+\pi f)M}\). We seek trial imaginary densities that
convert this oscillation into a positive real cost while preserving
the prescribed totals. The estimates here are variational statements;
the particle contour deformation implementing them is proved in
Section~\ref{sec:04-gas-fluctuations}.

Let \(a=\operatorname{Re}\eta,\ b_1=\operatorname{Im}\eta\) (here \(b_1\) is not a size); suppose \(|a|\le a_0<\pi,\ |b_1|\le b_0\) fixed. Allow Fourier aliases \(j\in\mathbb Z\). In the real negative-log action \(\mathcal A\), add to \(\operatorname{Re}\mathcal E\) the negative log modulus of \(e^{i(\eta+\pi f+2\pi j)M}\). Use \(D_z=\partial^{-1}\sum z_\sigma\), \(z_r=z_+\), so \(\int z_r=-\operatorname{Im}M\). Let \(f_R^\psi\) be the cumulative positive \emph{weighted} smeared far count to the right, and \(f_R^0\) the unsmeared version. The sign tails and \eqref{eq:G-1} give
\begin{equation}
\begin{split}
 \mathcal A\ \ge\
 B(x)&-B(z)-2p\int(D z_r+r D_z)-2\sum_{I:\rm far}g_I\int x_I D_z
 -(x,J_o z)\\
 &-\int \Theta z_r-b_1\int r-s(h)(B(x)+B(z)+1)-C_h n_e,
\end{split}
 \label{eq:G-4}
\end{equation}
Here \(\Theta=a+2\pi(j+f_R^\psi)\). Indeed the baseline \(+\) versus \(-\) sign tails give the \(2p\) terms by total neutrality and antisymmetry. The additional far rows give sign against \(z_++z_-\) with coefficient \(g_I\), and against \(z_r\) with coefficient \(\pi w_I\); the constant from \(2 f_R^\psi-f\) thus cancels the sign-bias \(\pi f\).

The next proposition starts with the unshifted real action and then gives
three uses of an imaginary trial field.  Part 2 permits every Fourier
alias.  Part 3 suppresses aliases whose effective phase is large near an
end of the wall.  Part 4 gives the quadratic dependence on the tilt when
those end phases are bounded.

We make the common assumptions and the end phases precise.  An admissible
near density \(\alpha=(\alpha_+,\alpha_-)\) is supported on \([-L,L]\),
has the prescribed totals \(m'_\sigma(M)\), and has finite trial energy.
We allow signed real parts.  Put
\[
 x=\operatorname{Re}\alpha-\mu+\psi*\zeta,
 \qquad z=\operatorname{Im}\alpha,
 \qquad
 \mathcal A=\operatorname{Re}\mathcal E(\alpha,\zeta)
             -\log\left|e^{i(\eta+\pi f+2\pi j)M}\right|.
\]
Thus \(\int z_\sigma=-\sigma\operatorname{Im}M\).
For parts 2--4 below assume that \(p_L\) has uniformly bounded
\(H^{-1/2}\) norm.  For each sufficiently small fixed \(\delta>0\),
assume that it has a normalized smooth truncation \(p_L^\circ\), supported
in \(|u|\le(1-\delta/2)L\), whose difference from \(p_L\) is exponentially
small in total variation and \(H^{-1/2}\).  The derivatives of each fixed
order of \(p_L^\circ\) are to be polynomially bounded in \(L\).
The hard-wall coordinates constructed below supply these properties.

For a fixed \(\delta\), choose nonnegative smooth probes
\(\varphi_\pm\le1\), supported away from the ends of
\(\pm[(1-2\delta)L,(1-\delta)L]\), obtained by rescaling fixed bumps,
with \(\int\varphi_\pm\asymp\delta L\).  For
\(a=\operatorname{Re}\eta\), define
\[
 \Theta(u)=a+2\pi\bigl(j+f_R^\psi(u)\bigr),\qquad
 \theta_\pm=\frac{\int\Theta\varphi_\pm}{\int\varphi_\pm},
 \qquad U=\max\{|\theta_+|,|\theta_-|\}.
\]
Here \(f_R^\psi\) is the weighted smeared far count to the right,
as in \eqref{eq:G-4}.  A polynomial bound throughout the proposition
means a fixed power of \(L\) with a fixed multiplicative constant.

\begin{proposition}[Trial bounds with tilt]\label{prop:03-trial-bounds}
Fix \(a_0<\pi\) and \(b_0<\infty\), and let
\(\eta=a+ib_1\) satisfy \(|a|\le a_0\), \(|b_1|\le b_0\).
The following bounds hold uniformly for admissible real parts of
\(\alpha\), with the imaginary fields specified below.  In parts 1--3,
\(h>0\) is fixed and sufficiently small.
\begin{enumerate}
\item If \(M\) is real and \(z=0\), then for sufficiently small
fixed \(h>0\), with no truncation of the explicit counts,
\[
 \mathcal A\ge cB(x)-C_hL-C_hn_e.
\]

\item Suppose \(n_e\) is polynomially bounded, and that \(p_L\)
has the properties stated above.  For any
\(v=\operatorname{Im}M\in\mathbb R\) and any \(j\in\mathbb Z\),
the choice \(z_\sigma=-\sigma vp_L\) gives
\[
 \mathcal A\ge cB(x)-C_hL-C_hn_e
       +2\pi jv-C_h\bigl(v^2+(1+f)|v|\bigr).
\]
There is no polynomial restriction on \(j\) in this estimate.

\item Suppose in addition that \(|j|\) is polynomially bounded.
There are a value \(v=\operatorname{Im}M\), independent of
\(\operatorname{Re}M\), and an imaginary field \(z\) with the
required totals such that
\[
 \mathcal A\ge cB(x)+c_\delta U^2L-C_hL-C_hn_e.
\]
The value \(v\) is polynomially bounded.  Each
\(z_\sigma+\sigma vp_L\) is the derivative of a smooth primitive
supported in \(|u|\le(1-\delta/2)L\), with polynomial bounds on
all required fixed-order derivatives.  Moreover, the number
of far particles in
\([-(1-2\delta)L,(1-2\delta)L]\) is at most \(CU\) for large \(L\).

\item Under the same polynomial hypotheses, suppose \(U\le U_0\).
Choose a desired error \(\varepsilon>0\), then \(h\) sufficiently
small, then \(\delta\) sufficiently small.  The constant \(U_0\)
may be fixed arbitrarily after these choices.  There are a value
\(v=\operatorname{Im}M\), independent of \(\operatorname{Re}M\),
and an imaginary field \(z\), with the same support and polynomial
regularity properties as in part 3, such that
\[
 \mathcal A\ge c'_\varepsilon B(x)
       +2L\frac{\operatorname{Re}(\eta^2)}{4\kappa}
       -C\varepsilon L-o(L)-C_hn_e.
\]
The selected \(v\) is polynomially bounded.
The retained coefficient \(c'_\varepsilon>0\) can be taken
independent of \(\delta,U_0\).  On these bounded-\(U\) configurations,
integrating the explicit factorial measures with the factor
\(\exp(C_hn_e)\), while using a fixed portion of the retained
\(B(x)\)-damping, costs at most
\[
 \exp\bigl(C_{h,\varepsilon}\delta L+o(L)\bigr).
\]
The same conclusion allows further fixed exponential factors per
explicit particle after reducing \(\delta\) as needed.
\end{enumerate}
The constants are uniform once the polynomial bounds and the previously
fixed parameters are fixed.  In part 4 the asymptotic limit is taken
only after \(\varepsilon,h,\delta,U_0\) have been chosen in that order;
constants depending on those fixed choices are allowed in \(o(L)\).
\end{proposition}
\begin{proof}
For part 1, set \(z=0\) in \eqref{eq:G-4}. The remaining
linear term satisfies
\[
 |b_1\textstyle\int r|\le C_h\sqrt{L(B(x)+1)}
                    \le \tfrac14B(x)+C_hL
\]
by Corollary~\ref{cor:03-counts}; choose \(h\) so that the smoothing
loss consumes at most another fixed fraction of \(B(x)\).
This proves the first bound without any count truncation.

For part 2 suppose that \(n_e\) is polynomially bounded and use the
stated norm and truncation properties of \(p_L\). For \(z_\sigma=-\sigma v p_L,\ v=\operatorname{Im}M\), \eqref{eq:G-4} gives
\begin{equation}
 \mathcal A\ge c B(x)-C_h L-C_h n_e+2\pi jv-C_h(v^2+(1+f)|v|).
 \label{eq:G-5}
\end{equation}
Here \(D_z=0\), and the mixed baseline and odd-remainder errors cost \(C|v|\sqrt{B(x)}\). This bound uses no polynomial restriction on \(j\).

For part 3 suppose also \(|j|\le{\rm poly}(L)\), with the polynomial fixed. We select \(v=\operatorname{Im}M\), independent of \(\operatorname{Re}M\), together with \(z\). In the choices below we first specify smooth \(z_r,D_z\) supported on \(|u|\le(1-\delta)L\) (and \(v=-\int z_r\)); then we actually add \(v(p_L^\circ-p_L)\) to \(z_r\), leaving \(D_z\) unchanged. Thus each difference \(z_\sigma+\sigma v p_L\) is a smooth density derivative with primitive compactly supported in \(|u|\le(1-\delta/2)L\); all its required sizes and fixed-order derivatives are polynomially bounded. In evaluating \eqref{eq:G-4} we can neglect the exponentially small correction, at cost \(o(1)(1+\sqrt{B(x)})\) by the norm bounds, total variation, and the polynomial count bounds.

Take smooth nonnegative probes \(\varphi_\pm(u)\le1\), supported on \(\pm[(1-2\delta)L,(1-\delta)L]\), rescaled fixed bumps supported away from their interval edges with mass comparable to \(\delta L\). Set \(\theta_\pm\) equal to the corresponding weighted averages of \(\Theta\); let \(U=\max|\theta_\pm|\). A first, coarse tilt uses
\[
 D_z=0,\qquad z_r=-t_0\sum_\pm\theta_\pm\varphi_\pm
\]
with small fixed \(t_0>0\). Then \(-\int\Theta z_r=t_0\sum\theta_\pm^2\int\varphi_\pm\). The \(B(z)\) cost is at most \(C t_0^2\sum\theta_\pm^2\int\varphi_\pm\); baseline and odd mixed terms are absorbed similarly with a fraction of \(B(x)\). The term with \(b_1\) is absorbed as in the real bound. Thus
\begin{equation}
 \mathcal A\ge cB(x)+c_\delta U^2 L-C_h L-C_h n_e .
 \label{eq:G-6}
\end{equation}
By positivity, the number of far particles on \([-(1-2\delta)L,(1-2\delta)L]\) is \(\le C U\) for large \(L\): their smeared masses contribute at least a constant each to the drop between the two probe averages.

For part 4 suppose \(U\le U_0\), where \(U_0\) may be arbitrarily large but fixed. Given a desired small error, choose a small parameter \(\varepsilon>0\), then \(h\) sufficiently small (so \(s(h)\ll\varepsilon\)), then \(\delta>0\) sufficiently small relative to these. \(U_0\) can be chosen after these parameters. Use a smooth mask \(0\le w(u)\le1\) supported on \([-(1-3\delta)L,(1-3\delta)L]\), vanishing at distances \(\le\sqrt L\) of far positions, equal to 1 except in boundary buffers of total width \(O(\delta L)\) inside \([-L,L]\) and holes of total width \(O_{U_0}(\sqrt L)\). Take fixed-order positive derivatives \(O_{U_0,\delta,l}(L^{-l/2})\). This is possible by the count conclusion. Put \(s(u)=\operatorname{sign}(a+2\pi(j+f_R^0(u)))\) and \(A_0=|a|/(2\kappa),\ B_0=-p b_1/(2\kappa)\). Use
\[
 z_r=-A_0\,s(u)w(u),\qquad D_z=B_0 w(u),
\]
with the correction convention above. These are smooth before correction even at jumps. Since \(|a|<\pi\), \(|a+2\pi(j+f_R^0)|\ge |a|\). On the mask support all far smearing tails from points outside their own respective side of \(u\) are superpolynomially small. Consequently \(-\int\Theta z_r\ge |a| A_0\int w-o(1)\), and \(\sum g_I\int x_I D_z=o(1)\). Moreover
\[
 B(z)\le (B_0^2+(\kappa-p^2)A_0^2)\int w+o(L),\qquad
 |(x,J_o z)|\le o(\sqrt L)\sqrt{B(x)} .
\]
Indeed before correction the scaled \(L^2\) masses of \(z_r,D_z\) (squared divided by \(L\)) concentrate at low frequency since their positive-order derivatives have norms \(o(\sqrt L)\). Use the symbol expansion and the upper norm bound, and the odd-remainder multiplier bound.

We give details for the real-field optimization so that no slow-variation assumption on \(x\) is needed. Let subscript \(\ell\) mean a smooth compact low-pass Fourier cutoff equal to 1 near zero, even real with values in \([0,1]\), supported on sufficiently small fixed frequencies depending on \(\varepsilon\). By the positive small-frequency expansion, reserving a fraction on all frequencies,
\[
 B(x)\ \ge\ c_\varepsilon B(x)
 +(1-C\varepsilon)\int \big(D_\ell^2+Q(r_\ell,x_{{\rm far},\ell})\big).
\]
On the mask \(x_{{\rm far},\ell}\) is superpolynomially small (Schwartz kernel tails, count bound). Thus in this lower bound we can use \((1-C'\varepsilon)\int w(D_\ell^2+(\kappa-p^2)r_\ell^2)-o(1)\) besides the reserve. The pairings against \(ws,w\) can replace \(D,r\) by \(D_\ell,r_\ell\) at \(o(\sqrt L)\sqrt{B(x)}\) cost by the derivative bounds and Sobolev norms. Also
\[
 |\int(1-w)r|\ \le\
 (C\sqrt{\delta L}+o(\sqrt L))\sqrt{B(x)}+C_h .
\]
Use the smooth wall plateau of \eqref{eq:G-3} minus \(w\), whose \(L^2\) norm has the stated bound by the support estimate and whose derivative norm is \(o(\sqrt L)\), then positive outer tail bounds for the rest. Constants depending on \(U_0,\delta,h\) in these estimates are allowed inside little-oh terms.

Absorb errors with a portion of the reserve, including \eqref{eq:G-1}, losing \(C\varepsilon L+o(L)\) with \(\delta\) sufficiently small. It remains on measure \(w\,du\) to complete squares for \(D_\ell,r_\ell\), whose linear coefficients in \eqref{eq:G-4} are \(2p A_0 s,\ -(b_1+2p B_0)\). The lower cost from these optimizations per unit mass is at worst
\[
 -p^2 A_0^2-\frac{(b_1+2p B_0)^2}{4(\kappa-p^2)}-O(\varepsilon).
\]
Adding \(-B_0^2-(\kappa-p^2)A_0^2+|a| A_0\) leaves \((a^2-b_1^2)/(4\kappa)\) before errors. Thus
\begin{equation}
 \mathcal A\ \ge\ c'_\varepsilon B(x)+2L\,\frac{\operatorname{Re}(\eta^2)}{4\kappa}
 -C\varepsilon L-o(L)-C_h n_e.
 \label{eq:G-7}
\end{equation}
The retained constant can be taken independent of \(\delta,U_0\). For these bounded-\(U\) estimates the explicit factorial measures with \(\exp(C_h n_e)\) and a portion of the retained damping cost at most \(\exp(C_{h,\varepsilon}\delta L+o(L))\): outer main points cost bounded grand mass by tail damping, only a bounded number of far points can be on the interior interval just counted, and the remaining far points use an effective volume \(O(\delta L)+O_{h,\varepsilon}(1)\) per variable, with factorial divisors. This also works with more fixed exponential factors per particle, by further reducing \(\delta\). Bounds \eqref{eq:G-5}–\eqref{eq:G-7} are deterministic bounds on trial densities. Implementing them for the complex integrals, including fluctuation determinants and the discreteness of \(M\), takes further work below.


\end{proof}

\subsection{Coordinates that remove the diagonal singularity}

The trial estimates require a probability coordinate concentrated near
\(\varrho\), while the reference ensemble uses a fixed compact
interval. We choose a map that has both properties and makes the
true-minus-reference interaction smooth through the endpoints.

Use a reference parameter \(s\in[-A,A]\), \(A=L+d_0\), and \(u=W(s)\). This is odd and \(W(A-v)=L-\omega(v)\) at the positive end, where \(\omega\) is quadratic near 0, Gevrey, increasing with derivative comparable to \(\min(v,1)\), and \(\omega(v)=v-d_0\) for \(v\ge1\). Construct it by integrating a cutoff interpolation of the derivative between \(v\) and 1 (fixes \(d_0\)). Thus \(W=s\) except within bounded endpoint distance and \(W\) reflects evenly at the ends. Set
\[
 \lambda_L(s)=C_L\varrho(W(s)),\qquad
 t=\theta(s)=\pi\int_{-A}^s\lambda_L,\qquad
 C_L^{-1}=\int_{-A}^{A}\varrho(W(v))\,dv=1+O(e^{-\gamma L}).
\]
So \(p_L\) is \(\lambda_L ds\) pushed to physical \(u\). Its stated probability and truncation norm requirements hold: in the core \(s=u\), and near each wall its density in \(u\) is bounded by \(C e^{-\gamma L}(L-|u|)^{-1/2}\); the log-kernel estimate for the \(H^{-1/2}\) norm thus suffices. Cut smoothly in the core to form \(p_L^\circ\). Derivatives of \(\lambda_L\) are bounded by \(C^{l+1}(l!)^{P}\lambda_L\) for a fixed \(P\), and \(\lambda_L\asymp e^{-\gamma|s|}\). Constants in Gevrey bounds below can change (including order \(P\)), but not with \(N\). These elementary closure estimates under smooth Gevrey multiplication, composition on a bounded scale, and reciprocal of a function bounded away from zero follow by multiplying formal Taylor series with coefficients bounded by \(C^l(l!)^{P}\); in composition of positive degrees totaling \(l\), products of factorials are bounded by \(l!\), and the number of compositions is at most exponential.

Let
\[
 K(s,v)=H_{nn}(W(s)-W(v))-\mathbf 1\,H_{\rm ref}(\theta(s),\theta(v)),
 \qquad \mathscr K=\partial_s\partial_v K ,
\]
where \(\mathbf 1\) denotes the identity matrix.

\begin{lemma}[Regularity of the compensated kernel]\label{lem:03-kernel-regularity}
The kernel \(K\) is smooth, extending evenly by reflection in each endpoint variable. \(K\) and all derivatives of total order \(l\) satisfy bounds \(C^{l+1}(l!)^P(1+L)\), and
\(\mathscr K\) and its derivatives satisfy \(C^{l+1}(l!)^P e^{-c|s-v|}\). For the latter estimates one only extends locally across endpoints. Here and below \(\operatorname{poly}(L)\) prefactors may replace \(1+L\) without harm.

\end{lemma}
\begin{proof}
First suppose \(|s-v|>c_0>0\). The sinh terms have the required estimates directly. Split the cos-difference logarithm into its two sine logarithms. Then
\(\sin(|\theta(s)-\theta(v)|/2)\) is bounded below by \(c\max_{[s,v]}\lambda_L\), reversing the interval if needed, and \(\sin((\theta(s)+\theta(v))/2)\) is at least as large. One derivative in each variable supplies the decaying ratio
\[
 \frac{\lambda_L(s)\lambda_L(v)}{(\max_{[s,v]}\lambda_L)^2}.
\]
Higher derivatives preserve this bound, proving the off-diagonal estimate for \(\mathscr K\).

Near the diagonal and away from the endpoints, divide both \(2\sinh((W(s)-W(v))/2)\) and \(\cos\theta(s)-\cos\theta(v)\) by \(s-v\). Their natural local scales are respectively \(1\) and \(\lambda_L^2\), since \(\sin\theta\asymp\lambda_L\) there. Averaging derivatives along the joining interval gives uniform Gevrey bounds and lower bounds on those scales. Their logarithmic ratio is therefore smooth.

Near a common endpoint, let \(d_s,d_v\) be the distances from that endpoint. Both differences have the additional factor \(d_s+d_v\): even reflection makes their zeros occur at both \(s=v\) and the reflected diagonal. Divide by
\[
 (s-v)(d_s+d_v).
\]
The division is smooth Gevrey, obtained by averaging a derivative of the first divided difference from its reflected zero. The quadratic fold \(\omega'(d)\asymp d\) and the sine factors give positive lower bounds for both quotients on their respective scales. The same log-ratio argument now proves regularity through the corner. Even reflection in either variable follows from the construction of \(W\) and \(\theta\).

Finally, in variables at distance at least \(\delta L\) from the endpoints, the compensated kernels continue analytically for uniformly small complex displacements, with the same fixed-order derivative control in the remaining variables. Here \(W(s)=s\) and \(\theta'=C_L\pi\varrho\) is analytic with relative bounds on uniformly small disks, so the divided-difference arguments apply there as well.
\end{proof}

The next statement supplies the linear and constant terms needed together with this kernel in the particle deformation and dimension interpolation. It explains why the exponentially large mean particle number does not make the compensated coefficients exponentially large.

\begin{lemma}[Regularity of the compensated source]\label{lem:03-source-regularity}
Suppose \(|M|,n_e,\ell_e\) have fixed polynomial bounds in \(L\). In the \(s\)-coordinate, the linear source
\[
 b=H_{nn}(m'p_L-\mu)+\widetilde H_{ne}\zeta
\]
in \(\Phi_M\) is smooth and even under endpoint reflection. Its derivatives of order \(r\) are bounded by
\(\operatorname{poly}(L)C^{r+1}(r!)^P\), for a fixed Gevrey order \(P\), and it has the core analytic continuation just described. Writing
\[
 E_e=\tfrac12(\zeta,H\zeta)_{\rm off},
\]
the difference \(\mathcal E(m'p_L,\zeta)-E_e\) is polynomially bounded in absolute value. The unchanged explicit-pair term \(E_e\), independent of \(M\), has real part bounded below by a negative polynomial; no upper bound is required.
\end{lemma}
\begin{proof}
For each main component, decompose the centered mean as
\[
 m'_\sigma p_L-m\varrho
   =(m'_\sigma C_L-m)\varrho
       +m'_\sigma(p_L-C_L\varrho).
\]
The coefficient in the first term is polynomially bounded, since \(m'_\sigma-m\) is polynomial and \(C_L-1=O(e^{-\gamma L})\). Its convolution is controlled by \(H_{nn}*\varrho=V\).

The two measures in the second term agree in the core. Near or beyond the ends their densities, written in \(ds\) for \(p_L\) and \(du\) for \(\varrho\), have derivative bounds with a factor \(O(e^{-\gamma L})\) and exponentially bounded tails. This factor cancels \(m'_\sigma=O(e^{\gamma L})\). Insert a fixed-width physical Gevrey cutoff there. For convolution of the first density with the diagonal logarithm, subtract \(-2\log|s-v|\) and its reflected-image term using the divided differences of \(W\). The remaining kernel is smooth. Transfer derivatives of the ordinary translated logarithm to the compactly cut off, evenly reflected density. For convolution with the second density, transfer derivatives in physical \(u\) and then compose with \(W\). Both terms are analytic in the deep core by separation from their supports.

The source \(\widetilde H_{ne}\zeta\) has the same bounds for fixed \(h\). Its regularization concerns only exterior main interactions and leaves an analytic kernel in the core. These arguments also give polynomial total variation and moments for \(m'p_L-\mu\), hence the stated constant-term bound by log integrability. Finally, real logarithmic singularities of explicit pairs are repulsive; away from them the kernels have at most linear growth. The polynomial count and excursion bounds therefore give \(\operatorname{Re}E_e\ge-\operatorname{poly}(L)\), including when distinct explicit particles approach each other.
\end{proof}

\subsection{The relative interaction operator}

The compensated kernel is now smooth. It remains to compare its
quadratic fluctuations with those of the reference law, uniformly as
the interval grows. Let \(Q=(Q_+,Q_-)\) be a primitive of a change
of zero mass on each main label in the \(s\)-coordinate, initially
smooth and zero at both endpoints. Define
\[
 h_0(Q,Q)=(Q',H_{\rm ref}Q'),\qquad
 H_{\rm form}(Q,Q)=h_0(Q,Q)+
   \iint Q(s)^{\mathsf t}\mathscr K(s,v)Q(v)\,ds\,dv
\]
where \(H_{\rm form}\) is the true near interaction form. The second expression is complex symmetric; on complex fields both energy forms are extended bilinearly. The space \(\mathcal X\) has squared norm \(h_0\); in the \(t\)-coordinate the fields extend as odd functions at both ends, completed from smooth ones. Per label
\[
h_0=\pi^2\sum_{n\ge1} n |b_n|^2,\qquad Q=\sum b_n\sin(nt)
\]
for the real form/norm, since \(H_{\rm ref}(t,t')=4\sum_{n\ge1}\cos(nt)\cos(nt')/n\) as an integrable-kernel Fourier expansion and \(\int\cos(nt)\,dQ=n\pi b_n/2\) on Dirichlet fields. For symmetric energy calculations on complex fields \(h_0(\cdot,\cdot)\) extends bilinearly; norms and operator adjoints use the ordinary complex Hilbert convention. The inverse kernel using \(ds\) pairing is
\[
 G(s,v)=\frac{1}{2\pi^2}\log\left(\frac{\sin((\theta(s)+\theta(v))/2)}
                  {\sin(|\theta(s)-\theta(v)|/2)}\right).
\]
In particular \(j:\mathcal X\to L^2(ds)\) the inclusion has bounded norm uniformly, \(jj^*=G\) separately per label: the log singularity is square integrably bounded locally by \(C(1+|\log|s-v||)\) using \(\sin\theta\le C\lambda_L\); off diagonal the kernel and derivatives decay exponentially, by writing the squared ratio as \(1+\sin\theta(s)\sin\theta(v)/\sin^2((\theta(s)-\theta(v))/2)\). Schur's integral bound applies.

The next estimate controls the number of nearly degenerate directions,
not just the smallest eigenvalue. This distinction is needed when
passing from trace estimates to logarithmic determinants.

\begin{lemma}[Trace bound and small eigenvalues]\label{lem:03-spectral}\label{E:spectral-bounds}
On the complexification of \(\mathcal X\), let
\[
 T=I+j^*\mathscr K j,
 \qquad A_T=\tfrac12(T+T^*).
\]
There are constants independent of \(L\) such that
\(\|T-I\|_1\le CL\), \(\|T\|\le C\), and
\begin{equation}
 A_T\ge cL^{-2}I,
 \qquad N_{A_T}(t):=\dim\mathbf1_{(-\infty,t)}(A_T)
       \le CL\sqrt t\quad(0<t<t_0).
 \label{eq:G-8}
\end{equation}
The same lower gap and counting estimate hold for the singular values of
\(T\), and for every orthogonal compression of \(T\).
\end{lemma}
\begin{proof}
Partition \([-A,A]\) into intervals \(I_a\) of lengths in \([1,2]\),
and use normalized cosine bases \(e_{a,n}\), \(n\ge0\), on these intervals.
The mixed-kernel estimates imply
\[
 \left|\big\langle e_{a,n},\mathscr K e_{b,m}\big\rangle\right|
 \le \frac{Ce^{-c|a-b|}}{(1+n^2)(1+m^2)}.
\]
For a positive frequency, integrate twice by parts: the first boundary
term vanishes because the sine is zero at either endpoint, and the
second boundary term involves only the first derivative of the kernel.
Doing this in both variables gives the displayed estimate, including
all derivative boundary terms. Summing the trace norms of these
rank-one matrix entries gives
\(\|\mathscr K\|_1\le CL\). Schur's test applied directly to its
exponentially decaying kernel gives \(\|\mathscr K\|\le C\).
The bounds for \(T\) follow from the uniform bound for \(j\).

It remains to control the lower edge of the spectrum. For a two-component
field, sum over components in the following definition:
\[
 \mathfrak e(Q)=
 \iint_{|s-v|\le2}\frac{|Q(s)-Q(v)|^2}{|s-v|^2}\,ds\,dv
 +\int_{d(s)\le2}\frac{|Q(s)|^2}{d(s)}\,ds,
 \qquad d(s)=A-|s|.
\]
We first establish the two comparisons
\[
 h_0(Q,Q)\le C\bigl(\mathfrak e(Q)+\|Q\|_{L^2(ds)}^2\bigr),
 \qquad \langle Q,A_TQ\rangle_{\mathcal X}\ge c\mathfrak e(Q).
\]
The circle difference formula for the odd extension in \(t\) expresses
\(h_0\), up to fixed constants, using the two kernels
\[
 \frac{\theta'(s)\theta'(v)}
      {\sin^2((\theta(s)-\theta(v))/2)},\qquad
 \frac{\theta'(s)\theta'(v)}
      {\sin^2((\theta(s)+\theta(v))/2)},
\]
paired respectively with \(|Q(s)-Q(v)|^2\) and \(|Q(s)+Q(v)|^2\).
When \(|s-v|\ge c_0\), the first kernel is exponentially integrable
in the separation. The second kernel has the same property except at
a common endpoint; there it is bounded by a constant times the inverse
square of the sum of the endpoint distances. Its integral in one
variable is at most \(C/d(s)\). These assertions follow from
\(\lambda_L(s)\asymp e^{-\gamma|s|}\) and
\(\sin\theta(s)\asymp\lambda_L(s)\min(1,d(s))\).
On the remaining diagonal region the first kernel is bounded by
\(C|s-v|^{-2}\). This proves the upper comparison.

For the lower comparison, push the primitive forward to physical
coordinate \(u=W(s)\) and extend it by zero outside \([-L,L]\).
Lemma~\ref{lem:03-coercivity} bounds its Fourier seminorm with multiplier
\(k^2/(1+|k|)\). For any fixed sufficiently large \(C_0\), this multiplier
is comparable to the multiplier
\[
 \int_{|z|\le C_0}\frac{|e^{ikz}-1|^2}{z^2}\,dz
 \asymp\frac{k^2}{1+|k|}
\]
of the truncated difference seminorm
\[
 \iint_{|u-u'|\le C_0}
 \frac{|Q(W^{-1}(u))-Q(W^{-1}(u'))|^2}{|u-u'|^2}\,du\,du',
\]
where the zero extension is understood. For an interior point at distance \(\ell<C_0/2\) from a wall,
integration against the adjacent exterior interval gives weight
\(1/\ell-1/C_0\asymp1/\ell\). Thus interactions with the exterior
zeros give a wall weight comparable to
\(W'(s)/(L-|W(s)|)\asymp1/d(s)\) near an endpoint.
When the two endpoint distances are comparable,
\(W'(s)W'(v)/|W(s)-W(v)|^2\asymp|s-v|^{-2}\).
The part with noncomparable distances is bounded by the wall term:
for example, integrate \(|s-v|^{-2}\) over \(d(v)<d(s)/2\) or
\(d(v)>2d(s)\), using \(|Q(s)-Q(v)|^2\le2|Q(s)|^2+2|Q(v)|^2\).
A fixed enlargement of the endpoint region causes no change in these
bounds. This proves the second comparison.

For completeness, divide any interval of length comparable to
\(b\ge1\) into cells of lengths in \([1/2,1]\), and let \(q_a\)
be the mean of \(Q\) on cell \(a\). The local seminorm controls
\[
 \sum_a\int_{I_a}|Q-q_a|^2
       +\sum_a|q_{a+1}-q_a|^2.
\]
The first bound is the elementary variance identity on a cell; the
second follows by averaging the difference over two adjacent cells,
whose mutual distances are at most \(2\). If the mean on the whole
block vanishes, the discrete Poincare inequality gives
\(\sum_a|q_a|^2\le Cb^2\sum_a|q_{a+1}-q_a|^2\).
Thus \(\|Q\|_2^2\le Cb^2\mathfrak e(Q)\) for fields with zero
mean on each such block. On the whole interval, the wall term controls
the first and last cell means; summing the successive differences then
gives the same estimate with \(b\asymp L\), without any constraints.
The two comparisons now give the gap \(cL^{-2}\).

For \(t\) above a fixed multiple of \(L^{-2}\), choose
\(b=c_1t^{-1/2}\) with \(c_1\) sufficiently small. Vanishing block
means impose at most \(CL/b\) linear constraints, and on their common
kernel the Rayleigh quotient of \(A_T\) is at least \(t\).
The min--max principle proves the counting bound. Below the global
gap the count is zero; the intervening range is handled by changing
the constants. Compressing a positive Hermitian form cannot increase
this count. Finally, if a subspace satisfies
\(\|TQ\|<t\|Q\|\), then on it
\(\langle Q,A_TQ\rangle\le\|Q\|\|TQ\|<t\|Q\|^2\).
The variational characterization of singular values therefore gives
the same count and gap for \(T\) and its compressions.
\end{proof}

\subsection{The determinant contribution}

Away from the center and the endpoints, the compensated kernels
approach translation-invariant kernels. Their Fourier symbol determines
the leading determinant. Lemma~\ref{lem:03-spectral} supplies the
uniform control near zero needed to integrate its logarithm.

\begin{theorem}[Relative determinant asymptotics]\label{thm:03-determinant}\label{E:determinant}
The Fredholm determinant satisfies
\begin{equation}
 -\tfrac12\log|\det T|=(7\gamma/8)L+o(L).
 \label{eq:G-9}
\end{equation}
\end{theorem}
\begin{proof}
Accretivity gives an injective operator with closed range. Since
\(T-I\) is trace class, \(T\) has Fredholm index zero and is
therefore invertible. Its Fredholm determinant is defined by the convergent exterior-power series.
Equivalently, take determinants of finite orthogonal compressions
converging strongly to the identity, then pass to the limit in trace
norm. The same procedure gives
\(\log|\det T|=\tfrac12\operatorname{Tr}\log(T^*T)\).
We compute this trace through the limiting symbol
\[
 T_\infty(k)=
 \frac{k^2\widehat H_{nn}(k)}
      {2\pi k\coth(\pi k/(2\gamma))}.
\]

First consider a nonconstant word in \(T-I\) and \(T^*-I\).
By cyclicity its trace is an integral alternating kernels
\(G\) and \(\mathscr K\) or \(\mathscr K^*\).
The local bound \(C(1+|\log|s-v||)\) for \(G\), and exponential
off-diagonal bounds for all kernels, provide an integrable majorant
for the relative coordinates in this cyclic integral. More explicitly,
truncate every separation at a fixed radius \(R\). The discarded
integral is bounded by \(CL e^{-cR}\) times a fixed polynomial in
\(R\). Each local logarithmic singularity can be integrated first
against an adjacent bounded \(\mathscr K\) kernel. With \(R\) fixed,
the central region and the two endpoint regions of width \(O(R)\)
contribute \(O(R)\), which vanishes after division by \(2L\).

Under translation into either remaining tail, with distance from both
the center and the endpoints tending to infinity, one has \(W(s)=s\)
and \(\theta(s)\) or \(\pi-\theta(s)\) asymptotic to a constant times
\(e^{-\gamma|s|}\), locally with every fixed derivative. Consequently
\[
 G(s,v)\longrightarrow
       \frac{1}{2\pi^2}\log\coth\frac{\gamma|s-v|}{2}.
\]
The mixed derivative of the reference kernel tends to that of
\(-2\log|2\sinh(\gamma(s-v))|\); its additional one-variable terms
disappear upon taking mixed derivatives. The Fourier transforms of
these two limiting kernels are reciprocal, respectively
\[
 \frac{\tanh(\pi k/(2\gamma))}{2\pi k},
 \qquad 2\pi k\coth\frac{\pi k}{2\gamma}.
\]
It follows from the logarithmic transform identities that multiplication
by the limiting \(\mathscr K\) symbol gives \(T_\infty-I\).
Dominated convergence, followed by \(R\to\infty\), proves that the
trace of each word, divided by \(2L\), converges to the integral of the
corresponding matrix trace against \(dk/(2\pi)\).

We pass from words to the logarithm without losing uniformity at zero.
For singular values \(s_i(T)\), Equation~\eqref{eq:G-8} gives
\[
 \sum_{s_i<\tau}\log\frac{\tau}{s_i}
 =\int_0^\tau\frac{\#\{s_i<t\}}{t}\,dt
 \le CL\sqrt\tau.
\]
Thus clipping \(\log(T^*T)\) below \(\tau^2\) changes its normalized
trace by at most \(C\sqrt\tau\). On the compact interval containing
its spectrum, approximate the clipped logarithm by polynomials
vanishing at one, with error at most \(\epsilon|y-1|\): apply polynomial
approximation to the continuous quotient by \(y-1\).
The resulting trace error is bounded by
\(\epsilon\|T^*T-I\|_1\le C\epsilon L\).
The symbol has the same integrable control: its smaller singular value
is comparable to \(k^2\) near zero, its other singular value stays
positive there, and it tends exponentially to the identity at infinity.
These assertions follow directly from the displayed \(2\times2\)
symbol and its positive Hermitian part. Hence clipping costs
\(O(\sqrt\tau)\) on the symbol side as well. Passing successively
\(L\to\infty\), \(\epsilon\downarrow0\), and \(\tau\downarrow0\)
yields
\[
 \lim_{L\to\infty}\frac{\log|\det T|}{2L}
       =\int_{\R}\log|\det T_\infty(k)|\,\frac{dk}{2\pi}.
\]

Set \(x=\beta|k|\) and \(y=e^{-x}\). The determinant is positive,
and its positive square root equals
\[
 \sqrt{\det T_\infty(k)}
 =\frac{c_*(k)\sinh(3x)}{\sinh(4x)\coth(3x/2)}
 =\frac{(1-y^3)(1-y^6)}{(1+y)(1-y^8)}.
\]
Integrating the absolutely integrable logarithmic series gives
\[
 \begin{split}
 \int_0^\infty\log\sqrt{\det T_\infty(x/\beta)}\,dx
 &=-\frac{\pi^2}{18}-\frac{\pi^2}{36}
       -\frac{\pi^2}{12}+\frac{\pi^2}{48}\\
 &=-\frac{7\pi^2}{48}.
 \end{split}
\]
Here \(\int_0^\infty\log(1-e^{-ax})\,dx=-\pi^2/(6a)\)
and \(\int_0^\infty\log(1+e^{-x})\,dx=\pi^2/12\), obtained by
termwise integration; \(\sum n^{-2}=\pi^2/6\) follows, for example,
from Parseval's identity. Accounting for both signs of \(k\), the
square of the displayed square root, and \(\gamma=\pi/(3\beta)\),
we obtain
\[
 \log|\det T|=-\frac{7\pi}{12\beta}L+o(L)
             =-\frac{7\gamma}{4}L+o(L),
\]
which is Equation~\eqref{eq:G-9}.
\end{proof}

\subsection{A finite-dimensional core for particle transport}

A determinant calculation on the completed Hilbert space does not yet
provide a deformation of a finite particle integral. We next approximate
the relative interaction by smooth fields supported where the mean
particle density is exponentially large. The remaining directions have
a negligible normalized logarithmic cost. Fixing the approximation
tolerance before taking the large-size limit is essential.

\begin{lemma}[Smooth core approximation]\label{lem:03-core}\label{E:finite-core}
For every sufficiently small fixed \(\xi>0\) and all sufficiently
large \(L\), there are \(\delta_*>0\), independent of \(L\), and a
real orthogonal projection \(P\) in \(\mathcal X\), of rank \(O_\xi(L)\),
whose range consists of smooth fields supported in
\(|s|\le(1-\delta_*)L\), such that
\begin{equation}
 \|T-I-P(T-I)P\|_1\le\xi L.
 \label{eq:G-10}
\end{equation}
An orthonormal basis of its range can be chosen with every required
fixed derivative bounded by a polynomial in \(L\). The images of these
basis fields under \(h_0\) and \(H_{\rm form}\), regarded as test
densities for the \(ds\) pairing, have the same bounds and smooth odd
reflections at the endpoints. Constants may depend on \(\xi\) and on
the derivative order.
\end{lemma}
\begin{proof}
We first approximate \(\mathscr K\) on \(L^2(ds)\). Truncate the cosine
basis in each unit-scale cell at frequency \(J\). The coefficient
estimate in Lemma~\ref{lem:03-spectral} gives trace-norm error at most
\(CL/J\), since \(\sum_{n>J}(1+n^2)^{-1}\le C/J\) and the sum over
cell separations converges. Approximate the finitely many retained
cosines in each cell by smooth functions compactly supported inside
that cell, with \(L^2\)-error at most \(\eta\), and orthonormalize.
For fixed \(J\) and sufficiently small \(\eta\), their Gram matrices
are uniformly invertible. Replacing a rank-one term \(u\otimes v\)
by its smoothed counterpart changes its trace norm by at most
\(\|u-\widetilde u\|\|v\|+
\|\widetilde u\|\|v-\widetilde v\|\).
Summing the absolutely convergent cell expansions therefore adds at
most \(C_J\eta L\).

Discard cells meeting \(|s|>(1-2\delta_*)L\). There are
\(O(\delta_*L+1)\) such cells, and the total trace-norm contribution
of terms with either index in these cells is of that same order.
Denote the resulting real bump projection by \(Q_b\). Then
\[
 \|\mathscr K-Q_b\mathscr KQ_b\|_1
 \le C\bigl(J^{-1}+C_J\eta+\delta_*\bigr)L+O(1).
\]
Choose \(J\), then \(\eta\), then \(\delta_*\), so the coefficient
on the right is as small as required. There are \(O_\xi(L)\) bumps;
each has uniformly bounded derivatives of every fixed order, with
constants allowed to depend on these fixed choices.

Let \(b_1,\ldots,b_d\) be the orthonormal bump basis and put
\(v_i=j^*b_i\). As a function in \(s\), \(v_i=Gb_i\).
Their Gram matrix in \(\mathcal X\) is uniformly conditioned. The
upper bound follows from \(\|j\|\le C\). For the lower bound, every
\(b\) in the bump span is a permissible smooth field and satisfies
\(\|b\|_{\mathcal X}\le C_\xi\|b\|_2\), by the local upper norm
estimate and its uniformly bounded \(H^1/L^2\) ratio. Duality gives
\[
 \|j^*b\|_{\mathcal X}
 \ge\frac{|\langle j^*b,b\rangle_{\mathcal X}|}
          {\|b\|_{\mathcal X}}
 =\frac{\|b\|_2^2}{\|b\|_{\mathcal X}}
 \ge c_\xi\|b\|_2.
\]
Thus the map \(b\mapsto Gb\) loses no linear independence or uniform
conditioning.

Choose a smooth cutoff equal to one on a neighborhood of the bump
support and zero outside \(|s|\le(1-\delta_*)L\), leaving a separation
at least \(c\delta_*L\) between the bumps and its nonconstant region.
Multiply each \(v_i\) by this cutoff, obtaining \(\widetilde v_i\).
The off-diagonal derivative bounds for \(G\) imply
\[
 \|v_i-\widetilde v_i\|_{\mathcal X}
       \le C_\xi L^C e^{-c\delta_*L}.
\]
To check the endpoint term, the explicit sine ratio gives
\[
 G(s,v)=O\bigl(d(s)e^{-c|s-v|}\bigr)
\]
when \(s\) is close to an endpoint and \(v\) lies in a retained cell.
Every fixed positive-order derivative has the same exponential bound,
without necessarily retaining the factor \(d(s)\). The difference
field vanishes at the endpoint. In particular, the elementary estimate
\[
 \int_0^2\frac{|q(d)|^2}{d}\,dd
 \le 2\int_0^2|q'(d)|^2\,dd,\qquad q(0)=0,
\]
controls its wall contribution.
The upper local norm estimate, or its consequence
\(h_0(q,q)\le C\|q\|_{H^1(ds)}^2\) for such fields, then proves the
claim. The Gram matrix remains uniformly invertible because
\(d=O_\xi(L)\). If \(P_0,P\) are the orthogonal projections onto the
spans of \(v_i,\widetilde v_i\), respectively, the formula
\(V(V^*V)^{-1}V^*\) for a span projection yields
\[
 \|P-P_0\|\le C_\xi L^C e^{-c\delta_*L}.
\]
The operator \(B=j^*Q_b\mathscr KQ_bj\) has both its range and its
adjoint range in \(\operatorname{ran}P_0\); hence \(B=P_0BP_0\).
Writing \(A_1=T-I\), we obtain
\[
 \|A_1-PA_1P\|_1
 \le2\|A_1-B\|_1+2\|P-P_0\|\,\|B\|_1
 \le C\bigl(J^{-1}+C_J\eta+\delta_*\bigr)L+o(1).
\]
This proves \eqref{eq:G-10} after the parameter choices above.

Finally, \(Gb_i\) has controlled derivatives of any fixed order in
the core. Locally subtract the ordinary diagonal logarithm from \(G\);
the remaining divided differences are smooth with uniform derivative
bounds. For the logarithm itself, transfer derivatives to the smooth
bump, and use the exponentially decaying off-diagonal estimates on the
rest. Multiplication by the cutoff preserves these bounds. Passing to
an orthonormal basis uses the inverse square root of a uniformly
conditioned \(d\)-dimensional Gram matrix; its operator norm is bounded,
so pointwise derivative bounds increase by at most \(O(\sqrt d)\).
For the action of \(h_0\), write the test density as the derivative of
\(H_{\rm ref}\) applied to the derivative of the core field. Subtract
its local logarithm and transfer derivatives to that smooth input.
The even endpoint reflections of the kernel give odd reflections of
the resulting test density. Adding the smooth operator \(\mathscr K\)
gives the same assertions for \(H_{\rm form}\).
\end{proof}

The finite core carries almost all of the determinant. Eliminating
its coordinates also leaves a positive interaction on its orthogonal
complement. The following statement quantifies both assertions using
the small-value count, rather than the gap alone.

\begin{lemma}[Determinants and Schur complements]\label{lem:03-schur}\label{E:schur-comparison}
Let \(S=\operatorname{ran}P\), \(R=S^\perp\), and
\[
 T_e=T_{RR}-T_{RS}T_{SS}^{-1}T_{SR}.
\]
The Hermitian part \(\operatorname{Re}T_e=(T_e+T_e^*)/2\) has the gap
and small-eigenvalue count in \eqref{eq:G-8}. Moreover,
\begin{equation}
 \left|\log|\det T_{SS}|-\log|\det T|\right|=o_\xi(L),
 \qquad
 \sum_{\substack{\lambda<1\\
        \lambda\in\operatorname{spec}(\operatorname{Re}T_e)}}
       -\log\lambda=o_\xi(L).
 \label{eq:G-11}
\end{equation}
Eigenvalues are counted with multiplicity. Here the notation means
\[
 \lim_{\xi\downarrow0}\limsup_{L\to\infty}
       \frac{|o_\xi(L)|}{L}=0.
\]
It does not assert an \(o(L)\) error for a fixed \(\xi\). One may take errors bounded respectively by
\(C\xi^{1/5}L\) and
\(C\xi^{1/3}(1+|\log\xi|)L\), after the core approximation is made
with tolerance \(\xi\).
\end{lemma}
\begin{proof}
The gap for \(T_{SS}\) makes it invertible. Set
\[
 Jy=(-T_{SS}^{-1}T_{SR}y,y),\qquad y\in R.
\]
Then \(TJy=(0,T_ey)\), so
\[
 J^*(\operatorname{Re}T)J=\operatorname{Re}T_e,
 \qquad \|Jy\|\ge\|y\|.
\]
The gap follows immediately. If \(E\subset R\) is a subspace on which
the Rayleigh quotient of \(\operatorname{Re}T_e\) is less than \(t\),
the quotient of \(\operatorname{Re}T\) on \(JE\) is also less than
\(t\). Since \(J\) is injective, min--max proves the count. These
identities use Hermitian adjoints; the matrices of the bilinear energy
are, separately, complex symmetric in a real basis.

We make explicit how the counting bound controls logarithms. If the
positive eigenvalues or singular values of any operator under
consideration have counting function \(N(t)\le CL\sqrt t\), then
\[
 \sum_{s_i<\tau}\log\frac{\tau}{s_i}
     =\int_0^\tau\frac{N(t)}t\,dt\le2CL\sqrt\tau.
\]
The lower gap justifies the identity even for the finite-size operators.
Put \(B=T_{SS}\oplus I_R\). Equation~\eqref{eq:G-10} gives
\(\|T-B\|_1\le\xi L\), and \(\|T\|,\|B\|\le C\), whence
\(\|T^*T-B^*B\|_1\le C\xi L\).
The function \(y\mapsto\log\max(y,\tau^2)\) has Lipschitz constant
\(\tau^{-2}\). The trace difference of this function on two positive
operators is bounded by that constant times their trace-norm difference;
this follows first for finite matrices by the eigenvalue variation
inequality and then by finite compressions. Removing the clipping costs
\(CL\sqrt\tau\) by the preceding integral identity. Thus
\[
 \left|\log|\det T|-\log|\det B|\right|
       \le CL\bigl(\sqrt\tau+\xi\tau^{-2}\bigr).
\]
Take \(\tau=\xi^{2/5}\) to obtain the first assertion.

For the second assertion choose \(0<\alpha<t_0\) and split the
singular-value expansion of \(T_{SS}^{-1}\) as \(B_{\rm hi}+B_{\rm lo}\),
where
\[
 \operatorname{rank}B_{\rm hi}\le CL\sqrt\alpha,
 \qquad \|B_{\rm lo}\|\le\alpha^{-1}.
\]
This is precisely the count for singular values of \(T_{SS}\) below
\(\alpha\). Every block of \(T-B\) has trace norm at most \(\xi L\),
while the off-diagonal blocks have operator norm at most \(C\).
It follows that
\[
 \operatorname{Re}T_e=I+F+E,\qquad
 \operatorname{rank}F\le2CL\sqrt\alpha,
 \qquad \|E\|_1\le C\xi L(1+\alpha^{-1}),
\]
with \(F,E\) Hermitian. The factor two in the rank bound accounts for
taking a Hermitian part. Since \(\operatorname{Re}T_e>0\), interlacing
and integration of its counting function below one imply
\[
 \operatorname{Tr}(I-\operatorname{Re}T_e)_+
       \le\operatorname{rank}F+\|E\|_1.
\]
Indeed a rank-\(r\) perturbation changes a count by at most \(r\), and
integration over \((0,1)\) then bounds the contribution of \(F\) by
\(r\); the contribution of \(E\) is bounded by its trace norm.
For \(0<\tau<t_0\), the clipped negative logarithm satisfies
\[
 -\log\max(\lambda,\tau)
 \le\frac{\log(1/\tau)}{1-\tau}(1-\lambda)_+.
\]
Combining this bound with the logarithmic tail integral gives
\[
 \sum_{\lambda<1}-\log\lambda
 \le2CL\sqrt\tau+
   \frac{\log(1/\tau)}{1-\tau}
   \left(CL\sqrt\alpha+C\xi L(1+\alpha^{-1})\right).
\]
Taking \(\alpha=\tau=\xi^{2/3}\) proves the stated estimate. The same
argument with \(\|\operatorname{Re}T-I\|_1\le CL\) and fixed \(\tau\)
also gives \(\sum_{\lambda<1}-\log\lambda=O(L)\) for
\(\operatorname{Re}T\).
\end{proof}


The centered interaction has now supplied three pieces of information:
a positive cost for explicit particles, the leading interior determinant,
and a smooth finite core that captures that determinant up to a
vanishing normalized loss. The next section supplies the remaining
probabilistic input: bounds for linear and compensated quadratic
statistics under the normalized reference ensemble.
